% Sección a cargo de: Gilces Carranza José Ignacio
\section{Ficha del proyecto y visión}\label{sec:ficha}

El documento sigue el proceso de ingeniería de requisitos de la norma ISO/IEC/IEEE 29148:2018~\cite{iso29148}, que exige identificar a las partes interesadas y delimitar el sistema antes de definir requisitos. La Tabla~\ref{tab:ficha} resume los datos de identificación del proyecto.

\subsection{Ficha del proyecto}

\begin{table}[H]
\centering
\caption{Ficha del Proyecto Fin de Curso.}
\label{tab:ficha}
\small
\begin{tabularx}{\textwidth}{@{}L{4cm}X@{}}
\toprule
\textbf{Campo} & \textbf{Detalle} \\
\midrule
Nombre del sistema & \sistema{} para Establecimientos de Suministros Eléctricos en Quevedo. \\
Organización cliente & \organizacion{}, establecimiento comercializador de materiales eléctricos ubicado en el cantón Quevedo, provincia de Los Ríos. \\
Tipo de sistema & Sistema de información transaccional de ventas e inventario, con un módulo de recomendación de productos asistido por inteligencia artificial. \\
Problema & En horas pico la facturación en mostrador es lenta, y el personal sin formación electrotécnica no siempre puede recomendar un producto compatible o un sustituto seguro cuando el artículo solicitado está agotado. Esto alarga la fila de atención y provoca ventas perdidas. \\
Oportunidad & El conocimiento de los vendedores con experiencia sobre equivalencias, sustitutos y venta cruzada puede registrarse en el catálogo y ponerse a disposición de todo el personal en el momento de la venta. \\
Objetivo general & Especificar los requisitos de una plataforma de ventas para \organizacion{} que agilice el cobro en mostrador y apoye al personal en la recomendación de productos eléctricos compatibles y disponibles en inventario. \\
Equipo & Arteaga Alava Danela Dayana, Cedeño Avila Winston Damian, Gilces Carranza José Ignacio (líder), Mendoza Palma Allan Jeremy. \\
\bottomrule
\end{tabularx}
\end{table}

\subsection{Declaración de visión}

Para redactar la visión se usa la plantilla de palabras clave que Wiegers y Beatty recomiendan para la declaración de visión del producto~\cite{wiegers2013}:

\begin{quote}
\textit{Para} los cajeros, administradores y propietarios de establecimientos de suministros eléctricos en Quevedo, \textit{que} enfrentan demoras en la facturación en horas pico y pérdida de ventas por no poder recomendar componentes compatibles o sustitutos, \textit{la} \sistemacorto{} \textit{es un} sistema de gestión comercial \textit{que} agiliza el cobro en mostrador, emite los comprobantes electrónicos exigidos por el SRI y sugiere productos eléctricos equivalentes que tienen existencias reales en el inventario. \textit{A diferencia de} la consulta manual de catálogos impresos o de los sistemas de facturación que solo registran ventas, \textit{nuestro producto} une el control de inventario con la asistencia técnica en el momento de la venta, sin ofrecer artículos agotados.
\end{quote}

% Pendiente (viernes): objetivos de negocio medibles (mínimo 3), con indicador,
% línea base obtenida en la organización y meta.
